% ============================================================================
\clearpage
\section*{v4.9.0-r01 Source-Only: The Photon Frame as Neutral Axis and the Conversion of Constants}
\addcontentsline{toc}{section}{v4.9.0-r01 Source-Only: The Photon Frame as Neutral Axis and the Conversion of Constants}
\label{sec:v490-r01-neutral-axis-conversion}
% ============================================================================

\paragraph{Purpose.}
Finding~(a) of \S\ref{subsec:v480-r18-findings} records that the conversion of constants between the matter and the
photon frame uses \(\gL=\cosh2\eta\), the matter--tachyon cross term, where the matter--photon pair is a neighbour pair
with factor \(\cosh\eta\).  This layer records where the factor came from, the working assumption under which it was
chosen, and what the assumption implies for a conversion law over three frames.  The computations are
\texttt{llm/k37\_gl\_umrechnungen.py}.  No value is compared with data.  Earlier layers are not rewritten.

\noindent\textbf{Source-only status.}
This v4.9.0-r01 layer is source-only.  It does not build a PDF, does not assign a new DOI,
and does not authorize a public release.  The v4.9.0 Version DOI
\href{https://doi.org/10.5281/zenodo.23534954}{10.5281/zenodo.23534954} and the Concept DOI
\href{https://doi.org/10.5281/zenodo.19210728}{10.5281/zenodo.19210728} are unchanged by it.

\subsection*{v4.9.0-r01 provenance of the factor}
\label{subsec:v490-r01-provenance}

In the archived versions v2.28.0 to v2.29.0 the splitting factor reads
``\(\gL=\gBI/\gstring\approx1.197\)'', while the same versions already carry \(1/\bsym\approx1.820\), that is
\(\bsym\approx0.5495\).  The value \(1.197\) is \((1-\bsym^2)^{-1/2}\), the Lorentz factor of the matter-frame
offset; the ratio \(\gBI/\gstring\) is \(1.864\).  In v2.29.1 the value was changed to \(1.864\) and the change recorded
as an arithmetic correction (``previous version used incorrect \(\gL\approx1.197\)''); the note carried in
Corollary~\ref{cor:gue-frame} calls the earlier value \(\gBI\), which it is not (\(\gBI\approx0.2375\)).  Since then
every conversion uses \(1.864\).

\paragraph{Author's statement (2026-10-21).}
The photon frame was assumed to be neutral, and the work was meant to proceed on that neutral axis first.  The archive
records the change of value only as an arithmetic correction and does not mention a neutral axis.

\paragraph{Theoremlet v4.9.0-r01.1.}
The two values the factor has carried are the two cross terms of Theoremlet r18.1:
\[
1.197=\cosh\eta=\sqrt{\tfrac{1+\gL}{2}}\quad\text{(matter--photon)},\qquad
1.864=\cosh2\eta=\gL\quad\text{(matter--tachyon)},
\]
with \(\eta=\operatorname{artanh}\bsym\).  Both follow from the single input \(\gBI/\gstring\).

\paragraph{Proof sketch.}
\(\cosh2\eta=2\cosh^2\eta-1\) and \(\cosh\eta=(1-\tanh^2\eta)^{-1/2}\).  \(\square\)

\subsection*{v4.9.0-r01 the neutral assumption in Reading~B}
\label{subsec:v490-r01-neutral}

The assumption has structural support in the complex-vector reading.  The photon is the sector fixed by the whole duality
circle (\S\ref{subsec:v480-r07-duality-angle}), the branch point of the duality holonomy (Theoremlet r19.2), and the
centre of the real \(s\)-axis on which matter and tachyon are mirror images (\S\ref{subsec:v480-r19-photon-centred}).
On the rapidity axis of Theorem~\ref{thm:symmetric-frame} it sits at rapidity \(0\), between matter at \(+\eta\) and
tachyon at \(-\eta\).  If the photon frame is neutral, a value on the neutral axis is reached from the matter frame by the
neighbour step, and the factor carried before v2.29.1 is the factor of that step.

\paragraph{Theoremlet v4.9.0-r01.2.}
A conversion law \(X^{(b)}=\gamma_{ab}^{-k}X^{(a)}\) with \(\gamma_{ab}\) the Lorentz factor of the pair is not
consistent over the three frames: from matter to tachyon through the photon it gives
\((\cosh\eta)^{2k}=\bigl(\tfrac{1+\gL}{2}\bigr)^k\), directly it gives \(\gL^{\,k}\) (for \(k=1\): \(1.432\) against
\(1.864\); for \(k=6\): \(8.63\) against \(42.0\)).  A law with the Doppler factor, \(X^{(b)}=e^{-k\eta_{ab}}X^{(a)}\),
composes exactly.

\paragraph{Proof sketch.}
\(\cosh^2\eta\neq\cosh2\eta\) for \(\eta\neq0\), while \(e^{\eta}e^{\eta}=e^{2\eta}\).  \(\square\)

\paragraph{Consequence.}
With the Lorentz-factor law one pair at a time is consistent.  Working on one axis first, as the author describes, is the
regime in which the law of \S\ref{sec:phase-constants} is free of contradiction; the obstruction of r01.2 appears only
once both pairs are used.  Under the neutral assumption the primary pair is matter--photon, with factor \(\cosh\eta\);
\(\gL\) is the factor across the neutral axis, from matter to tachyon.

\subsection*{v4.9.0-r01 a frame-order finding in the Hawking proposition}
\label{subsec:v490-r01-hawking}

Items~(i) and~(ii) of Proposition~\ref{prop:hawking-gBI-tachyon} give \(T_H^{(\mathrm{photon})}=\gL\,T_H^{(\mathrm{matter})}\)
and \(T_H^{(\mathrm{tachyon})}=\gL^{-1}T_H^{(\mathrm{matter})}\).  Reciprocal factors place photon and tachyon on opposite
sides of matter.  On the axis of Theorem~\ref{thm:symmetric-frame} the photon lies between matter and tachyon.  A law along
the axis moves the tachyon value beyond the photon value in the same direction; a mirror law about a neutral photon gives
the tachyon the matter value.  Neither gives reciprocal factors, so item~(ii) carries a frame order that does not match the
axis.  This is also the one statement of the main body in which the same factor serves both pairs.

\subsection*{v4.9.0-r01 ways forward, recorded and not chosen}
\label{subsec:v490-r01-ways}

\begin{enumerate}[leftmargin=2em]
\item[(A)] Keep \(\gL\) and the Lorentz-factor law, and read each conversion of \S\ref{sec:phase-constants} as one across
the neutral axis (matter--tachyon).  Statements whose endpoint is the photon frame then need the neighbour factor.
\item[(B)] Use the neighbour factor \(\cosh\eta=1.197\) for matter--photon, the value carried before v2.29.1, in line with
the neutral assumption.  This changes the photon-frame numbers of the main body (table below).
\item[(C)] Use a composable law with the Doppler factor, \(e^{\eta}=\sqrt{(1+\bsym)/(1-\bsym)}=1.854\) per neighbour step.
It is consistent over three frames.  At \(\bsym\) its neighbour factor lies \(0.55\%\) below \(\gL\), so the numbers of the
main body change little while their meaning changes.  By Theoremlet r01.3 below the closeness is structural: in the
continuum limit of Theorem~\ref{thm:tribonacci} the two factors are equal.
\end{enumerate}

\paragraph{Theoremlet v4.9.0-r01.3.}
The full-span Lorentz factor equals the neighbour-step Doppler factor, \(\cosh2\eta=e^{\eta}\), exactly when \(e^{\eta}=1\)
or \(e^{\eta}=\mathbb T\), the tribonacci constant of Theorem~\ref{thm:tribonacci}.  At that point
\(\cosh2\eta=e^{\eta}=\coth\eta=\mathbb T=1.83929\), \(\beta=0.54369\): the self-referential condition
\(\gL=1/\bsym(\gL)\) of Theorem~\ref{thm:tribonacci} is the condition that the matter--tachyon Lorentz factor and the
matter--photon Doppler factor coincide.

\paragraph{Proof sketch.}
With \(x=e^{\eta}\), \(\cosh2\eta=(x^2+x^{-2})/2\), so \(\cosh2\eta=x\) is \(x^4-2x^3+1=0\), which factors as
\((x-1)(x^3-x^2-x-1)=0\).  Further \(\coth\eta=(x^2+1)/(x^2-1)\), and \(\coth\eta=x\) is again \(x^3-x^2-x-1=0\).
Checked numerically to \(10^{-16}\).  \(\square\)

\paragraph{Consequence for the neutral axis.}
In the continuum limit the author's choice of \(\gL\) for the step from matter to the neutral photon frame is exactly the
factor a composable law assigns to that step: reading (A) and reading (C) agree there.  With the measured
\(\gL=\gBI/\gstring\), which lies \(1.37\%\) above \(\mathbb T\) (the lattice correction of
Remark~\ref{rem:op1-tribonacci}), they differ by \(0.55\%\).  The Lorentz-factor law of reading (B) does not share this
property.

\begin{center}
\small
\begin{tabular}{@{}lrrr@{}}
\toprule
quantity in the main body & \(\gL\) & neighbour, Lorentz & neighbour, Doppler\\
\midrule
\(\lP^{(\mathrm{matter})}/\lP^{(\mathrm{photon})}\) (Theorem~\ref{thm:planck-photon}), power \(3/2\) & \(2.546\) & \(1.309\) & \(2.525\)\\
vacuum-energy ratio, power \(6\) & \(42.0\) & \(2.94\) & \(40.6\)\\
\(m_\nu^{(\mathrm{matter})}/m_\nu^{(\mathrm{photon})}\) (Proposition~\ref{prop:k-neutrino}), power \(4\) & \(12.08\) & \(2.05\) & \(11.82\)\\
matter error (Corollary~\ref{cor:matter-error}), power \(5\) & \(22.53\) & \(2.46\) & \(21.92\)\\
\bottomrule
\end{tabular}
\end{center}
The table is arithmetic on the formulas of the main body; photon-frame values are not measured.  Statements that are
compared with data, among them Corollary~\ref{cor:falsifiable} and Proposition~\ref{prop:muon-g2}, are not re-evaluated
here; a re-evaluation under (A), (B) or (C) is a prediction and requires a pre-registration note.

\subsection*{v4.9.0-r01 register entries}
\label{subsec:v490-r01-register}

\begin{center}
\small
\begin{tabular}{@{}>{\raggedright\arraybackslash}p{0.34\textwidth}>{\raggedright\arraybackslash}p{0.58\textwidth}@{}}
\toprule
statement & status after v4.9.0-r01\\
\midrule
conversion law of \S\ref{sec:phase-constants} & factor \(\gL\) is the matter--tachyon term; consistent for one pair at a time (r01.2); under the neutral photon assumption the matter--photon factor is \(\cosh\eta\); choice among (A)--(C) open\\
Cor.~\ref{cor:gue-frame}, v2.29.1 note & the earlier value \(1.197\) was \(\cosh\eta\), the matter--photon Lorentz factor, not \(\gBI\) (r01.1)\\
Prop.~\ref{prop:hawking-gBI-tachyon} (ii) & reciprocal factors put the tachyon on the far side of matter; does not match the axis order (r01)\\
Thm.~\ref{thm:tribonacci} & equivalent to \(\cosh2\eta=e^{\eta}\): matter--tachyon Lorentz factor equals matter--photon Doppler factor (r01.3)\\
\bottomrule
\end{tabular}
\end{center}

\subsection*{v4.9.0-r01 ledger and guard update}
\label{subsec:v490-r01-ledger-guard}

\[
R_{\rm conv}^{v490\text{-}r01}=(0,0,1)
\]
for (provenance of the factor recorded, consistency of the conversion law over three frames decided and its relation to
the tribonacci fixed point established, choice among (A)--(C) open).

The layer spends no guard:
\[
R_{\rm glob}^{(v490\text{-}r01,+)}=(0,1,1,1,1),
\qquad
\Gamma_{\rm guard}^{v490\text{-}r01}=(1,1,1,1,1),
\qquad
\beta_4=1/256\ \text{unspent}.
\]
It does not imply exact finite \(C5\) zero closure, exact all-readout terminality,
\(C6/T2\), a global photon-clock/RH theorem, or global all-readout closure.

